\section{Trace-Length Mechanism}
\label{sec:trace-length}
\label{sec:length}
\label{app:trace-length-optimization}


\begin{lemma}[Connected open-cover budget]
\label{lem:open-cover-budget}
Let $X$ be compact and connected, with a finite Borel measure $\mu$
positive on every nonempty relatively open subset. Let $U_1,\ldots,U_m$
be open sets covering $X$, put $T_j=\overline{U_j}$, and assume at least
two of the relative pieces $U_j\cap X$ are nonempty. Then
\[
 \mu(X)<\sum_j\mu(U_j\cap X)\le\sum_j\mu(T_j\cap X).
\]
Consequently closed-piece upper bounds whose sum is at most $\mu(X)$
already contradict the open cover.
\end{lemma}
\begin{proof}
If the nonempty relative pieces were pairwise disjoint, one and the union
of the others would disconnect $X$. Some pair therefore overlaps in a
nonempty relatively open set, of positive measure. Integrating the covering
multiplicity gives strict inequality. Closure only enlarges each piece.
\end{proof}

For $X=\partial H$ or $S$, use length; for $X=H$, use area. Each measure
has the stated full-support property. The six V roles meet each target
in positive measure; the C role does so on $H$ and $S$, but need not meet
$\partial H$. In particular,
\begin{equation}
 \Haus(E)<L_E(T_C)+\sum_{i=0}^5L_E(T_i),\qquad E\in\{\partial H,S\},
 \label{eq:length-subadditivity}
\end{equation}
whenever these targets are covered. Empty center boundary traces are simply
omitted from the collection. The local geometric caps remain separate from
this common final budget argument.

\begin{figure}[H]
\centering
\begin{tikzpicture}[scale=.84]
  \foreach \x/\title in {0/$\partial H$: length $6$,4/$D$: length $6$,8/$S=\partial H\cup D$: length $12$}{
    \begin{scope}[xshift=\x cm]
      \coordinate (a0) at (1,0);
      \coordinate (a1) at (.5,.866);
      \coordinate (a2) at (-.5,.866);
      \coordinate (a3) at (-1,0);
      \coordinate (a4) at (-.5,-.866);
      \coordinate (a5) at (.5,-.866);
      \coordinate (o) at (0,0);
      \draw[hex] (a0)--(a1)--(a2)--(a3)--(a4)--(a5)--cycle;
      \node[panellabel,align=center] at (0,-1.38) {\title};
    \end{scope}
  }
  \draw[demand] (1,0)--(.5,.866)--(-.5,.866)--(-1,0)--(-.5,-.866)--(.5,-.866)--cycle;
  \foreach \p in {(4,0),(4.5,.866),(3.5,.866),(3,0),(3.5,-.866),(4.5,-.866)}{
    \draw[radialdemand] (4,0)--\p;
  }
  \draw[demand] (9,0)--(8.5,.866)--(7.5,.866)--(7,0)--(7.5,-.866)--(8.5,-.866)--cycle;
  \foreach \p in {(9,0),(8.5,.866),(7.5,.866),(7,0),(7.5,-.866),(8.5,-.866)}{
    \draw[radialdemand] (8,0)--\p;
  }
  \node[figlabel,align=center,text=DemandRed] at (8,-1.93)
    {the full-skeleton estimate is used only under its stated branch hypotheses};
\end{tikzpicture}

\caption{The perimeter and full-skeleton targets.  The skeleton is used only
under the hypotheses of the corresponding trace budget.}
\label{fig:strategy1-targets}
\end{figure}

\subsection{Skeleton bounds}

\begin{lemma}[Center skeleton cap]
\label{lem:center-skeleton-cap}
\label{lem:app-center-skeleton-cap-calculation}

If $N_E(T_C)\in\{1,2\}$ and
\[
 \left|T_C\cap\{M_0,\ldots,M_5\}\right|=1,
\]
then
\[
 L_S(T_C)<\frac32.
\]
\end{lemma}
\begin{proof}
Normalize the unique midpoint as $M_0$ and a positive boundary trace to
$e_{0,1}$.  Use the signed parameters and six radial exits of
\zcref{prop:signed-center-normal-form}. The positive trace condition is
\begin{equation}
 \alpha+W\delta<P.
 \label{eq:center-skeleton-domain}
\end{equation}
Summing the boundary and radial contributions gives
\[
 L_S(T_C)\le K_R+\mathcal E(\alpha,\delta),
 \qquad K_R:=W+E-\frac\eta R,
\]
where
\[
 \mathcal E(\alpha,\delta)
 :=-\frac\alpha R+\frac\alpha W+\delta
 +\min\left\{\frac\alpha R,\frac\delta W\right\}
 +\frac{[P-(R\alpha+\delta)]_+}{W}.
\]
Splitting according to the minimum and positive part gives
\begin{equation}
 \mathcal E(\alpha,\delta)\le\frac PW.
 \label{eq:center-skeleton-E-bound}
\end{equation}
If $\alpha/R\le\delta/W$, the expression without the positive part is
$\alpha/W+\delta$; when that part is active, the excess over $P/W$ is
$\alpha-R\delta/W\le0$.  The reflected order gives
$\delta-W\alpha/R\le0$.  If the positive part vanishes,
\eqref{eq:center-skeleton-domain} gives the same bound strictly.

Put $A:=1+R-R^2$.  Then
\[
 \frac32-K_R-\frac PW=\frac{1+A-2EA}{2RW}.
\]
Both sides of $1+A>2EA$ are positive, and
\[
 (1+A)^2-4A^2E^2=R^2W^2(5+4R-4R^2)>0.
\]
Therefore $K_R+P/W<3/2$, as required.
\end{proof}

\begin{lemma}[Positive-support skeleton cap]
\label{lem:positive-support-skeleton-cap}
Let $T_i$ be the closure of an original open V triangle.  If
\[
 \exists j\in\{i-1,i+1\}\quad\Haus(T_i\cap r_j)>0,
\]
then
\[
 L_S(T_i)<\frac32.
\]
\end{lemma}

\begin{proof}
Normalize to $i=0$ and translate the hexagon by $V_0$.  The five components
of $S$ on which a unit triangle containing $V_0$ can have positive trace
become the corresponding local components of the translated skeleton.  The
positive adjacent-radial trace becomes a positive boundary trace there.
\zcref[S]{prop:unique-center-midpoint} and
\zcref{lem:center-skeleton-cap} give the bound; every nonlocal component
is more than unit distance from $V_0$, apart from zero-measure endpoints.
\end{proof}

\begin{lemma}[No-support skeleton cap]
\label{lem:no-support-skeleton-cap}
If
\[
 A_i+B_i\le1,
\]
\[
 \Haus(T_i\cap r_{i-1})=0,
 \qquad \Haus(T_i\cap r_{i+1})=0,
\]
then
\[
 L_S(T_i)\le2.
\]
\end{lemma}

\begin{proof}
For $s>0$ and $j\notin\{i-1,i,i+1\}$,
$\|V_i-sV_j\|^2=1+s^2-2s\cos((j-i)\pi/3)>1$.
Thus only the incident boundary traces and own radial trace contribute;
their total is $A_i+B_i+C_i\le2$.
\end{proof}

For the remaining cap, reflect the local directions if necessary so that
$0\le A_i\le B_i\le1$.  The actual reach triple of a supercritical V
triangle lies in the selected cell
\begin{equation}
\begin{aligned}
 A_i^2+A_iB_i+B_i^2&\le1,\\
 C_i&\le\frac12,\\
 (A_i^2-1)C_i^2+(2A_iB_i^2+B_i)C_i+(B_i^4-B_i^2)&\le0,
\end{aligned}
\label{eq:supercritical-admissible-cell}
\end{equation}
where $C_i$ belongs to the connected component containing $0$ in the
corresponding one-dimensional fiber.  This is Cell $S$ of
\zcref{prop:exact-admissible-set}.  The selector is essential: the
quadratic is negative at $C_i=0$, positive at $C_i=1/2$, and opens downward.

\begin{lemma}[Supercritical skeleton cap]
\label{lem:supercritical-skeleton-cap}
\label{lem:app-supercritical-skeleton-calculation}

If $A_i+B_i>1$, then
\[
 L_S(T_i)<\frac32.
\]
\end{lemma}
\begin{proof}
The corner normal form, \zcref{prop:vd-corner-normal-form}, makes
$\Vdone$ and $\Vdtwo$ triangles nonsupercritical.  \zcref[S]{lem:t3-nonsupercritical} gives the same conclusion for $\Tlike$
triangles.  Hence a supercritical triangle is $\Vdzero$ and has no positive
adjacent-radial support.  Reflect so that $A_i\le B_i$.
Then
\[
 L_S(T_i)=A_i+B_i+C_i.
\]

Set
\[
 s:=A_i+B_i,\qquad \delta:=B_i-A_i,
 \qquad \gamma_0:=\frac32-s.
\]
The diameter constraint gives
\[
 1<s\le\frac2{\sqrt3},
 \qquad0\le\delta\le\sqrt{4-3s^2}.
\]
Let $P(\gamma)$ be the quadratic in
\eqref{eq:supercritical-admissible-cell}.  Substitution and multiplication by
$16$ yield
\[
\begin{aligned}
16P(\gamma_0)=Q(s,\delta)
&=\delta^4+8\delta^3s-6\delta^3
 +14\delta^2s^2-18\delta^2s+5\delta^2\\
&\quad-8\delta s^3+30\delta s^2-34\delta s+12\delta\\
&\quad+s^4-6s^3-19s^2+60s-36.
\end{aligned}
\]
Its derivative factors as
\[
 \frac{\partial Q}{\partial\delta}
 =2(2\delta+4s-3)
 \bigl(\delta^2+\delta(4s-3)+(s-1)(2-s)\bigr)\ge0.
\]
Therefore
\[
 Q(s,\delta)\ge Q(s,0)
 =(s-1)(s^3-5s^2-24s+36)>0.
\]
The cubic decreases on $[1,2/\sqrt3]$ and at the right endpoint equals
$88/3-136\sqrt3/9>0$.  Hence $P(\gamma_0)>0$.  The selected smaller-root
component in \eqref{eq:supercritical-admissible-cell} forces
$C_i<\gamma_0$, so $A_i+B_i+C_i<3/2$.
\end{proof}

\begin{lemma}[Positive-support rescuer]
\label{lem:positive-support-rescuer}
Assume that the open roles cover $S$, that $N_E(T_C)\in\{1,2\}$, and,
after symmetry,
\[
 T_C\cap\{M_0,\ldots,M_5\}=\{M_0\}.
\]
If $N_+\ge2$, then for a supercritical index $s\ne0$ there is
$j\in\{s-1,s+1\}$ such that
\[
 M_s\in U_j,
 \qquad \Haus(T_j\cap r_s)>0.
\]
\end{lemma}

\begin{proof}
Choose a supercritical $s\ne0$.  \zcref[S]{lem:self-midpoint} gives
$M_s\notin T_s$, and the displayed midpoint identity gives
$M_s\notin T_C$.  The open cover therefore puts $M_s$ in another $U_j$.
Diameter locality restricts $j$ to $s-1$ or $s+1$, and openness makes its
trace on $r_s$ have positive length.
\end{proof}

\begin{theorem}[Direct CE1/CE2 skeleton count]
\label{thm:common-skeleton-count}
If
\[
 N_E(T_C)\in\{1,2\},
\]
\[
 \left|T_C\cap\{M_0,\ldots,M_5\}\right|=1,
\]
and
\begin{equation}
 N_++N_{\mathrm{sp}}\ge3,
 \label{eq:common-skeleton-count}
\end{equation}
then the seven triangles do not cover $S$.
\end{theorem}

\begin{proof}
By \zcref{prop:vd-corner-normal-form,lem:t3-nonsupercritical}, every
positive-support role is nonsupercritical. Hence the two index sets are disjoint. With $k=N_++N_{\mathrm{sp}}$,
the preceding caps give
\[
 L_S(T_C)+\sum_iL_S(T_i)
 <\frac32+\frac32k+2(6-k)=\frac{27-k}{2}\le12,
\]
contrary to \eqref{eq:length-subadditivity}.
\end{proof}

\subsection{Boundary trace bounds}

The following bounds suffice for the perimeter exit.

\begin{theorem}[Boundary trace bounds]
\label{thm:boundary-trace-table}
For closures of the original open roles,
\[
 L_{\partial H}(T_C)=0\quad(\CEzero),\qquad
 L_{\partial H}(T_C)<\tfrac12\quad(\CEone\text{ or }\CEtwo),
\]
and
\begin{equation}
 L_{\partial H}(T_i)=A_i+B_i\le\frac2{\sqrt3}.
 \label{eq:vertex-boundary-locality}
\end{equation}
A nonsupercritical role contributes at most $1$; a Vd1 or Vd2 role
contributes less than $1/2$.
\end{theorem}
\begin{proof}
Boundary locality and the diameter bound give
$A_i^2+A_iB_i+B_i^2\le1$, hence
$3(A_i+B_i)^2/4\le1$. The remaining caps are
\zcref{cor:signed-center-boundary,prop:vd-corner-normal-form}.
\end{proof}

\begin{proposition}[Vd2 adjacent-midpoint branch]
\label{prop:ce2-vd2-midpoint-length}
Suppose the C triangle is CE1 or CE2 and $N_+=1$. If a Vd2 role contains
an adjacent midpoint, the perimeter is not covered.
\end{proposition}
\begin{proof}
By \zcref{lem:vd2-neighbor-midpoint-cap}, that role contributes strictly
less than $1/3$. The common boundary caps give
\[
 L_{\partial H}(T_C)+\sum_iL_{\partial H}(T_i)
 <\frac12+\frac2{\sqrt3}+\frac13+4<6.
\]
Apply \zcref{lem:open-cover-budget}.
\end{proof}
